\documentclass[11pt]{article}

\usepackage[margin=1in]{geometry}
\usepackage{amsmath, amssymb, amsthm}
\usepackage{bm}
\usepackage{mathtools}
\usepackage{enumitem}
\usepackage{xcolor}
\usepackage{hyperref}
\usepackage{listings}
\usepackage{booktabs}

\hypersetup{
    colorlinks=true,
    linkcolor=blue,
    urlcolor=blue,
    citecolor=blue
}

\definecolor{codegray}{gray}{0.95}
\lstset{
    backgroundcolor=\color{codegray},
    basicstyle=\ttfamily\small,
    frame=single,
    breaklines=true,
    columns=fullflexible
}

\title{Q-Learning Explained\\
\large Learning Values Instead of Directly Learning a Policy}
\author{}
\date{}

\begin{document}

\maketitle

\section{Goal of This Note}

Q-learning is a value-based reinforcement learning algorithm.
Instead of directly learning a policy $\pi(a \mid s)$, it learns how good each action is in each state:

\[
\boxed{
Q(s,a) = \text{expected future reward after taking action } a \text{ in state } s.
}
\]

Once we know $Q(s,a)$, choosing actions is simple:

\[
a^*(s) = \arg\max_a Q(s,a).
\]

\section{Policy Learning vs Value Learning}

In policy gradient, we learn:

\[
\pi_\theta(a \mid s).
\]

This directly tells us the probability of each action.

In Q-learning, we learn:

\[
Q(s,a).
\]

This tells us how valuable each action is.
The policy is then derived from the Q-values:

\[
\pi(s) = \arg\max_a Q(s,a).
\]

\section{The Bellman Idea}

The central idea is recursive.
The value of an action equals immediate reward plus future value:

\[
\boxed{
Q(s,a)
=
r + \gamma \max_{a'} Q(s',a').
}
\]

This equation says:

\[
\text{value now}
=
\text{reward now}
+
\text{discounted best value later}.
\]

\section{Temporal Difference Target}

When the agent takes action $a$ in state $s$, it observes:

\[
(s, a, r, s').
\]

The Q-learning target is:

\[
y
=
r + \gamma \max_{a'} Q(s',a').
\]

The prediction is:

\[
Q(s,a).
\]

The temporal difference error is:

\[
\delta
=
y - Q(s,a).
\]

The word temporal means across time.
The estimate at time $t$ is improved using the estimate at time $t+1$.

\section{Tabular Q-Learning Update}

If the state space is small, we can store Q-values in a table.
The update rule is:

\[
\boxed{
Q(s,a)
\leftarrow
Q(s,a)
+
\alpha
\left[
r + \gamma \max_{a'}Q(s',a') - Q(s,a)
\right].
}
\]

Here:

\begin{itemize}[leftmargin=2em]
    \item $\alpha$ is the learning rate.
    \item $\gamma$ is the discount factor.
    \item $r + \gamma \max_{a'}Q(s',a')$ is the target.
    \item $Q(s,a)$ is the current prediction.
\end{itemize}

\section{A Tiny Table Example}

Suppose:

\[
Q(s,a) = 2.0,
\qquad
r = 1.0,
\qquad
\gamma = 0.9,
\qquad
\max_{a'}Q(s',a') = 4.0.
\]

Then:

\[
y = 1.0 + 0.9 \cdot 4.0 = 4.6.
\]

The TD error is:

\[
\delta = 4.6 - 2.0 = 2.6.
\]

If $\alpha = 0.1$, then:

\[
Q(s,a) \leftarrow 2.0 + 0.1 \cdot 2.6 = 2.26.
\]

The value moved a little toward the target.

\section{Minimal Tabular Code}

\begin{lstlisting}[language=Python]
from collections import defaultdict
import random

Q = defaultdict(lambda: [0.0, 0.0])  # two actions

alpha = 0.1
gamma = 0.99
epsilon = 0.1

def choose_action(state):
    if random.random() < epsilon:
        return random.choice([0, 1])
    return max(range(2), key=lambda a: Q[state][a])

def update(state, action, reward, next_state, done):
    best_next = 0.0 if done else max(Q[next_state])
    target = reward + gamma * best_next
    error = target - Q[state][action]
    Q[state][action] += alpha * error
\end{lstlisting}

\section{Exploration: Epsilon-Greedy}

If the agent always chooses:

\[
\arg\max_a Q(s,a),
\]

it may never discover better actions.

Epsilon-greedy exploration solves this by choosing:

\[
\begin{cases}
\text{random action} & \text{with probability } \epsilon, \\
\text{best known action} & \text{with probability } 1-\epsilon.
\end{cases}
\]

At the beginning of training, $\epsilon$ is often high.
Later, it is reduced.

\begin{lstlisting}[language=Python]
epsilon = max(0.05, epsilon * 0.995)
\end{lstlisting}

\section{Why Q-Learning is Off-Policy}

Q-learning learns the value of the greedy policy:

\[
\max_{a'} Q(s',a').
\]

But the agent may collect data using epsilon-greedy exploration.

So the behavior policy is exploratory, but the learned target is greedy.
This is why Q-learning is called off-policy.

\section{Q-Learning with Neural Networks}

For large or continuous state spaces, a table is impossible.
CartPole has continuous observations:

\[
s =
\begin{bmatrix}
x \\
\dot{x} \\
\theta \\
\dot{\theta}
\end{bmatrix}.
\]

We cannot make a table for every possible real-valued state.
So we approximate:

\[
Q_\theta(s,a)
\]

with a neural network.

For CartPole, the network can map:

\[
\mathbb{R}^4 \rightarrow \mathbb{R}^2.
\]

The two outputs are:

\[
\begin{bmatrix}
Q(s,\text{left}) \\
Q(s,\text{right})
\end{bmatrix}.
\]

\section{DQN Loss}

Deep Q-Networks, or DQN, train a neural network to predict Q-values.

The prediction is:

\[
Q_\theta(s,a).
\]

The target is:

\[
y
=
r + \gamma \max_{a'}Q_{\theta^-}(s',a').
\]

The symbol $\theta^-$ means a target network.
It is a slowly updated copy of the Q-network.

The loss is:

\[
\boxed{
\mathcal{L}(\theta)
=
\left(
Q_\theta(s,a) - y
\right)^2.
}
\]

\section{Minimal DQN Target Code}

\begin{lstlisting}[language=Python]
import torch
import torch.nn.functional as F

q_values = q_net(states)                         # shape: [batch, actions]
chosen_q = q_values.gather(1, actions[:, None])  # Q(s, a)

with torch.no_grad():
    next_q = target_net(next_states).max(dim=1).values
    target = rewards + gamma * (1 - dones) * next_q

loss = F.mse_loss(chosen_q.squeeze(), target)

optimizer.zero_grad()
loss.backward()
optimizer.step()
\end{lstlisting}

\section{Experience Replay}

DQN usually stores transitions in a replay buffer:

\[
(s,a,r,s',d).
\]

Then it samples random minibatches from the buffer.
This helps because consecutive environment steps are highly correlated.

\begin{lstlisting}[language=Python]
replay_buffer.append((state, action, reward, next_state, done))
batch = random.sample(replay_buffer, batch_size)
\end{lstlisting}

Replay makes learning more stable and data efficient.

\section{Q-Learning vs REINFORCE}

\begin{center}
\begin{tabular}{lll}
\toprule
Question & REINFORCE & Q-Learning \\
\midrule
What is learned? & Policy probabilities & Action values \\
Core object & $\pi(a \mid s)$ & $Q(s,a)$ \\
Action choice & Sample from policy & Argmax Q-value \\
Update signal & Return $G_t$ & TD target \\
Wait until episode end? & Usually yes & No \\
On/off-policy & On-policy & Off-policy \\
\bottomrule
\end{tabular}
\end{center}

\section{Important Intuition}

Q-learning does not ask:

\[
\text{Which action did my policy sample, and how good was the episode?}
\]

It asks:

\[
\text{What should the value of this state-action pair be,
given the reward and the best future action?}
\]

So the basic learning signal is:

\[
\boxed{
\text{prediction error between current Q-value and Bellman target.}
}
\]

\section{Minimal Mental Model}

The shortest way to remember Q-learning is:

\[
\boxed{
Q(s,a) \leftarrow \text{reward now} + \text{best discounted value later}.
}
\]

Once $Q$ is accurate, the policy is easy:

\[
\boxed{
\text{choose the action with the largest Q-value.}
}
\]

\end{document}
